\documentclass[11pt]{article}
\usepackage[margin=1in]{geometry}
\usepackage{amsmath,amssymb,amsthm,hyperref}
\newcommand{\Tr}{\operatorname{Tr}}
\newcommand{\diag}{\operatorname{diag}}
\newcommand{\norm}[1]{\left\lVert#1\right\rVert}
\newtheorem{theorem}{Theorem}
\newtheorem{lemma}{Lemma}
\title{The universal positive-block norm inequality forces an essentially Hermitian off-diagonal block}
\author{Principia Math in Progress}
\date{12 September 2026}
\begin{document}
\maketitle

\section*{Task 1. Resolve the essentially Hermitian conjecture}
All matrices are complex and all norms below are operator norms. A matrix is
\emph{essentially Hermitian} if it has the form $\alpha I+\beta T$ with
$T=T^*$ and $\alpha,\beta\in\mathbb C$. Equivalently, its numerical range is
a line segment, possibly a point. We prove Conjecture 3.3 of Bourin--Lee
\cite{BL}, also Conjecture 7.3.3 on page 114 of Bourin's thesis \cite{B}.

\begin{theorem}
Suppose $X\in M_n(\mathbb C)$ satisfies
\begin{equation}\label{universal}
 \norm{\begin{pmatrix}A&X\\X^*&B\end{pmatrix}}
 \leq \norm{A+B}
 \quad\text{whenever }\begin{pmatrix}A&X\\X^*&B\end{pmatrix}\geq0.
\end{equation}
Then $X$ is essentially Hermitian. No invertibility or multiplicity
restriction on $X$ is necessary.
\end{theorem}

\begin{proof}
We first prove normality. For any $A>0$, the matrix
\[
 M_A=\begin{pmatrix}A&X\\X^*&X^*A^{-1}X\end{pmatrix}
 =\begin{pmatrix}A^{1/2}\\X^*A^{-1/2}\end{pmatrix}
   \begin{pmatrix}A^{1/2}&A^{-1/2}X\end{pmatrix}
\]
is positive semidefinite. Its nonzero eigenvalues are those of
$A+A^{-1/2}XX^*A^{-1/2}$. Thus \eqref{universal} implies
\begin{equation}\label{schur}
 \norm{A+A^{-1/2}XX^*A^{-1/2}}
 \leq\norm{A+X^*A^{-1}X}.
\end{equation}
Fix a unit vector $v$ and $0<\varepsilon<1$, and put
$P=vv^*$ and $H=P+(1-\varepsilon)(I-P)$. For $n>1$, $H$ has simple
largest eigenvalue $1$ with eigenvector $v$; for $n=1$ normality is automatic.
Apply \eqref{schur} with $A=tH$ and divide by $t$. With $u=t^{-2}$,
\[
 \lambda_{\max}(H+uH^{-1/2}XX^*H^{-1/2})
 \leq\lambda_{\max}(H+uX^*H^{-1}X).
\]
For a Hermitian $C$, the simple-eigenvalue first variation is
\begin{equation}\label{variation}
 \lim_{u\downarrow0}\frac{\lambda_{\max}(H+uC)-1}{u}=v^*Cv.
\end{equation}
For completeness, a unit maximizing vector $w_u$ converges to the line
$\mathbb Cv$ by the spectral gap of $H$. The Rayleigh principle bounds
the quotient in \eqref{variation} below by $v^*Cv$ and above by
$w_u^*Cw_u$, proving the formula.
Consequently
\[
 \norm{X^*v}^2\leq v^*X^*H^{-1}Xv.
\]
Letting $\varepsilon\downarrow0$ gives
$v^*XX^*v\leq v^*X^*Xv$ for every unit vector $v$.
The positive semidefinite matrix $X^*X-XX^*$ has trace zero, so it is zero.
Hence $X$ is normal.

We next record a necessary symmetry of extreme eigenvalues. For any
Hermitian $H$, define
\[
 K_X(H)=\begin{pmatrix}H&X\\X^*&-H\end{pmatrix},\qquad
 a=\lambda_{\max}(K_X(H)),\quad b=-\lambda_{\min}(K_X(H)).
\]
The trace is zero, so $a,b\geq0$. The matrix $bI_{2n}+K_X(H)$ is positive
semidefinite and its two diagonal blocks sum to $2bI_n$.
Condition \eqref{universal} therefore gives $b+a\leq2b$, or $a\leq b$.
Writing $J=\diag(I_n,-I_n)$, we have
$K_X(-H)=-J K_X(H)J$. Applying the same argument to $-H$ gives $b\leq a$.
Thus
\begin{equation}\label{symmetry}
 \lambda_{\max}(K_X(H))=-\lambda_{\min}(K_X(H))
 \qquad(H=H^*).
\end{equation}

By simultaneous unitary conjugation in \eqref{universal}, we may now
assume $X$ diagonal. Suppose three of its eigenvalues $z_1,z_2,z_3$
are noncollinear. Set
\[
 \Delta=\Im\bigl(\overline{z_2-z_1}(z_3-z_1)\bigr)\ne0,\qquad
 D=\diag(z_1,z_2,z_3),\qquad
 L=\begin{pmatrix}0&i&1\\-i&0&1\\1&1&0\end{pmatrix}.
\]
Choose $h>2\norm{X}/\sqrt3$ and take $H=hL\oplus0_{n-3}$.
Up to a permutation, $K_X(H)$ is the direct sum of
\[
 K=\begin{pmatrix}hL&D\\D^*&-hL\end{pmatrix}
\]
and matrices with eigenvalues $\pm|z_j|$, $j>3$.
The eigenvalues of $L$ are $0,\sqrt3,-\sqrt3$.
The off-diagonal perturbation of $\diag(hL,-hL)$ has norm
$\norm D\leq\norm X$. The Rayleigh principle therefore gives
\[
 \lambda_{\max}(K)\geq\sqrt3h-\norm X>\norm X,\qquad
 \lambda_{\min}(K)\leq-\sqrt3h+\norm X<-\norm X.
\]
Both extreme eigenvalues of the full matrix consequently come from $K$.
By \eqref{symmetry} they are $R$ and $-R$ for some $R>0$.

We show that $K$ cannot have any pair of opposite nonzero eigenvalues.
Direct trace identities give
\begin{equation}\label{traces}
 \Tr K=\Tr K^3=0,\qquad \Tr K^5=-20h^3\Delta.
\end{equation}
Here is an explicit verification of the only nontrivial identity. Write
$K=S+T$, with $S=\diag(hL,-hL)$ and
$T=\left(\begin{smallmatrix}0&D\\D^*&0\end{smallmatrix}\right)$.
Words with an odd number of $T$'s have trace zero. Since $DD^*=D^*D$,
the remaining contributions to $\Tr K^5$ reduce, by cyclicity, to
\[
 5h^3\Tr\bigl(L^2(D^*LD-DLD^*)\bigr).
\]
Indeed the two cyclic classes with two $T$'s each have five words;
$\Tr(S^3T^2)=0$, and the class $S^2TST$ is the expression displayed.
The terms with zero or four $T$'s cancel between the diagonal blocks.
Using
\[
 L^2=\begin{pmatrix}2&1&i\\1&2&-i\\-i&i&2\end{pmatrix},
\]
the remaining trace equals
\[
 4\Im(z_1\overline z_2-z_1\overline z_3+z_2\overline z_3)
 =-4\Delta.
\]
The cubic trace in \eqref{traces} follows in the same way from
$\Tr S^3=0$ and $\Tr(ST^2)=0$.

Newton's identities now show that the characteristic polynomial has form
\[
 p(\lambda)=\lambda^6+c_2\lambda^4+c_4\lambda^2
               +4h^3\Delta\lambda+c_6.
\]
In detail, the coefficients of $\lambda^5$ and $\lambda^3$ vanish
because the first and third power traces vanish, and the coefficient of
$\lambda$ is $-\Tr(K^5)/5=4h^3\Delta$.
But $p(R)=p(-R)=0$ would give $8h^3\Delta R=0$, a contradiction.
Thus no three eigenvalues of $X$ are noncollinear. The spectrum lies
on a line (the assertion is also immediate for at most two distinct
eigenvalues). Normality then gives $X=\alpha I+\beta T$ with $T=T^*$,
as required.
\end{proof}

If ``positive'' is interpreted as positive definite, the same conclusion
holds: add $\delta I$ to both diagonal blocks of each semidefinite test
and let $\delta\downarrow0$. The source's known width-zero inequality
gives the converse for essentially Hermitian $X$ \cite{BL,B}.

\paragraph{Task 2. Quantify stability.}
For $\norm X\leq1$, relate the distance from $X$ to the essentially
Hermitian matrices to the largest violation of \eqref{universal} over
positive test matrices with bounded norm, tracking dependence on $n$.

\paragraph{Task 3. Classify the Schatten exponents.}
For each $1<p<\infty$, characterize the matrices $X$ for which the
universal inequality holds with the Schatten $p$-norm; compare the known
normality characterization at $p=2$ with the operator-norm result here.

\paragraph{Task 4. Minimize the detecting test families.}
Determine which restrictions on the ranks of the diagonal blocks and
the full positive test matrix still suffice to detect every matrix that
is not essentially Hermitian.

\paragraph{Task 5. Examine infinite-dimensional operators.}
Determine the corresponding universal-inequality class for bounded
operators on an infinite-dimensional Hilbert space, where the
finite-dimensional trace argument forcing normality is unavailable.

\begin{thebibliography}{9}
\bibitem{BL} J.-C. Bourin and E.-Y. Lee,
\emph{Eigenvalue inequalities for positive block matrices with the
inradius of the numerical range}, International Journal of Mathematics
33 (2022), 2250009; \href{https://arxiv.org/abs/2111.15180}{arXiv:2111.15180},
Conjecture 3.3, p.~6.
\bibitem{B} J.-C. Bourin, \emph{A Journey into Matrix Analysis},
\href{https://arxiv.org/abs/2307.03064}{arXiv:2307.03064} (2023),
Conjecture 7.3.3, p.~114.
\end{thebibliography}
\end{document}
